\section{DISCUSSION}
\label{sec:discussion}

\subsection{RQ1: What Concurrent Speech Reports About Within-Person Presence}

The primary RQ1 finding is that, on participants held out from pipeline selection, the speech-derived IPQ total converged with the self-reported IPQ total (Section~\ref{subsec:rq1_convergence}), and that within a participant the pipeline identified which VR session had felt more present (Section~\ref{subsec:rq1_within}). What it did not reproduce is the survey's level: speech-derived totals were about a quarter of a point lower, a difference that is statistically detectable but smaller than the questionnaire's own measurement error, as the level-offset paragraph below takes up.

Two features of the within-person result carry the interpretation. First, the five cases in which the pipeline missed a participant's highest-presence game all involved the same pair of titles: those participants rated \textit{Half-Life: Alyx} above \textit{Subside} on the survey, whereas the pipeline ranked \textit{Subside} higher (Section~\ref{subsec:rq1_within}). The errors are therefore tied to specific content, not scattered at random. Second, the INV-side calibration, which shifts level, left every direction indicator unchanged (Section~\ref{subsec:rq1_convergence}). Direction and level are thus separable properties of the speech-derived score, and the paper keeps them apart for that reason.

The level offset deserves a second reading. The quarter-point gap is statistically detectable, but it is smaller than the standard error of measurement of the questionnaire itself, and the two totals were equivalent within a margin derived from that error under every rule that uses the full-sample reliability (Section~\ref{subsec:rq1_convergence}). The offset was not stable across samples, though: it had been reduced to $-0.05$ on the development participants and re-emerged at $-0.38$ on the held-out participants, so the anchor calibration generalized rank and variation but not absolute level. Two qualifications keep this from becoming a stronger claim than the data support. Equivalence holds for group means, not for individuals: the limits of agreement span more than a scale point and the concordance correlation is $.34$, and when both scores are placed on the percentile classes of two decades of IPQ studies~\cite{tran2024ipq}, only a quarter of participants fall in the same class, with the speech-derived class the lower one in nearly every mismatch. And the margin is one we derived and specified after seeing the data, so the equivalence analysis is exploratory. What the result does license is a narrower statement: the pipeline's total lands within the questionnaire's own resolution at the group level, so the offset is a matter of calibration and does not indicate that the two instruments measure different things.\footnote{INV1, the item the pipeline almost never scores, also has the weakest item--total correlation in the reference instrument (corrected $r=.13$).}
\textcolor{orange}{Can we make this paragraph more concise? too much micro-explanations? }

The subscale results tell a second-order story. Within participants, GP ($r=+.42$), SP ($r=+.29$), and REAL ($r=+.42$) converged with the survey, whereas INV was effectively zero and slightly negative ($r=-.06$). INV was also the least-covered subscale in speech (Section~\ref{subsec:sample}). The human-coder comparison (Section~\ref{subsec:rq1_reliability}) shows that this is not simply what participants left unsaid: coders found INV1 evidence in 8 of the 12 compared sessions and INV3 in 10, whereas the pipeline scored them in none and one. The cause is the evidence gate at aggregation (Section~\ref{sec:system}): INV1 and INV3 are scored from the absence of a signal, so their cells carry no sub-indicator marked as present, and the gate admits an item only on present evidence. This is a deliberate trade of coverage on absence-scored items for a stricter evidence standard, and it is why the paper reports the total as the primary indicator and treats INV as a documented coverage gap, not a scoring failure.
\textcolor{orange}{
change expressions like ``... it is why this paper reports ...'' --> we postulate ... 
... again need to make concise and easy to understand ... too much stat jargons overall ...
? }

\textcolor{orange}{
I guess it would be difficult to project the reason of RQ1 results ...?  because it has to do with the AI and associated implementations?
}

\textcolor{orange}{
Can we have a punchline statement(s) about RQ 1 in bold?  after all the verbose statements, I am not clear on what is the final stance on RQ1 ...}

\subsection{RQ2: Front-End Reach}

RQ2 evaluates the locally fine-tuned routing classifier on classification accuracy and pipeline reach, not on IPQ score accuracy: RQ1 and RQ2 are not compared on the same axis. On the held-out set, the deployed configuration reached a macro $F_1$ of $0.567$, with the keyword layer buying recall on rare categories at the cost of some precision. On the audit sessions, $289$ of $294$ item slots received at least one candidate row and $823$ of $843$ scored items passed the evidence gate; the losses were concentrated at the scoring stage, not at routing or at the gate. The retrained adapter was not adopted because it lost SP recall relative to the deployed configuration. Together, these numbers describe an on-device front end that reliably feeds the downstream scoring chain and does not appear to be the source of the INV coverage gap: INV recall on the held-out set was $0.650$, in the middle of the range across categories.

\textcolor{orange}{
I guess it would be difficult to project the reason of RQ2 results ...?  because it has to do with the AI and associated implementations?
}

\textcolor{orange}{Similar problems, too many stat numbers and possibly redundant ... but nevertheless the section itself is short ! Lets have a last short punchline summary in bold}


\subsection{RQ3: Whether TAP Perturbed Presence}

The TAP group scored higher than the non-TAP group on every measure in Table~\ref{tab:group}, including workload and cybersickness. Both raw and Holm-corrected $p$ values are reported: on the raw $p$ the IPQ total and INV cleared $.05$, whereas after Holm--Bonferroni correction across the seven measures no measure remained significant. We therefore read the RQ3 result as a preliminary null against the reactivity hypothesis: concurrent TAP did not measurably degrade presence in this sample. Three caveats attach. The uniform elevation across measures cannot be separated from a general response-intensity difference with these data. The effect was not uniform across games and was absent for The Lab. The lower confidence bound on the IPQ total ($g=0.791$, 95\% CI $[0.086,\ 1.495]$) sits only marginally above zero. A larger sample and, if feasible, a within-participant manipulation in which the same participant plays both with and without the think-aloud instruction would tighten the conclusion.

One further observation bears on the reactivity question, although it concerns the consistency of self-report and not its level. The internal consistency of the spatial-presence subscale differed sharply between groups on participant-level item means: $\alpha=.869$ in the non-TAP group but $.191$ in the TAP group, while the totals were comparable (Section~\ref{subsec:sample}). If this pattern replicates, it would suggest that concurrent verbalization perturbs presence self-report not by shifting the mean but by loosening the coherence of the spatial-presence responses, a form of reactivity that a mean comparison alone would miss. With 16 participants per group the contrast is fragile, and we record it as a hypothesis for the larger study, not as a finding.

\textbf{Thinking aloud did not cost presence: concurrent verbalization added no measurable workload or sickness, and the TAP group reported presence at least as high as the group that played in silence.}

\subsection{Implications for VR Evaluation Practice}

The direction result implies that speech-derived IPQ scoring can complement the post-session survey when the practical question is a within-person comparison across content (``did this session feel more present than that session for this player?'') and no questionnaire can be given between sessions. The level offset implies that the same scores should not currently be used as absolute presence scores against an external reference; the pipeline undercounts by about a quarter of a point on the total and by about half a point on INV, before calibration. The evidence gate implies that the score for a session is a statement about the items the player's speech actually supported, and that an INV score should not be reported for a session in which INV was uncovered; imputing zero would misrepresent the speech. In short, the current system is a within-participant comparator that documents its own coverage.

\paragraph{Cost relative to manual coding.}
The operating cost of the pipeline can be placed against what manual coding of the same corpus would require. Scoring the 3{,}595 TAP utterances by hand would mean judging each utterance against the 14 IPQ items with a $0$--$6$ rating; deductive coding of single-label documents has been timed at 72--144\,s per document~\cite{chew2023llm}, and our multi-item task is heavier, so we bracket the per-utterance time at 60--180\,s. With two independent coders and 10\% adjudication, as verbal-protocol analysis conventionally requires~\cite{chi1997quantifying}, this amounts to 132--396 person-hours. manual coding would cost roughly USD 1{,}350--4{,}050, before transcription, which adds 3--10 hours per hour of recording if done by hand~\cite{bailey2008transcribing}. The pipeline's recorded API expenditure for the scoring stages was USD 0.98--2.22 per session, or about USD 50--110 for all 48 sessions, and the sequential run completed in about two days of wall-clock time.\note{TODO: confirm wall-clock time from the run ledger; state the exchange rate date} These are operating costs: the development iterations that produced the ontology, guides, and prompts are not included, and the comparison says nothing about the quality of manual coding, which was not measured. Every intermediate judgment (routing, drafted evidence, review decision, rating) is logged, which makes the automated path auditable in a way that manual coding notes rarely are.
